\section{Normal-boundary continuation and transversality of the physical clock}
\label{sec:cc-clock}
A regular critical record is the actual $m$-flight record of
Section~\ref{sec:cr-critical} with $p_0=p_m=0$.  It is one half of a
true reversing periodic trajectory.  We vary the common physical
disk radius within $[R_-,R_+]$, keeping the ordered lattice centers
$O_0,\ldots,O_m$, all $m$ collisions and the unit-speed clock.  A
formal reflection continuation will always be distinguished from a
physical continuation.  The latter is required to avoid every third
disk and preserve the actual first-hit order.

\subsection{The full normal-boundary reflection equations}
Choose local real lifts of the contact angles $\alpha_i$ and of the
flight directions $\vartheta_i$.  Their normal-reflection equations
are, with the constants fixed at the reference record,
\begin{equation}\label{eq:cc-reflection}
 E_0=\alpha_0-\vartheta_0,\quad
 E_i=2\alpha_i-\vartheta_{i-1}-\vartheta_i\ (1\leq i<m),\quad
 E_m=\alpha_m-\vartheta_{m-1}.
\end{equation}
At a physical critical record $E_0$ is zero modulo $2\pi$ and
$E_m$ is $\pi$ modulo $2\pi$.  The interior constants are the
appropriate fixed odd multiples of $\pi$.  These lifts remove any
branch-cut ambiguity in the following derivatives.

\begin{theorem}[A full reflection inverse without a grazing divisor]
\label{thm:cc-full-resolvent}
Let $\mathcal B=D_\alpha E$ at a regular critical record with $m\geq1$.
Write $w_i=\tau_i^{-1}$ and set $w_{-1}=w_m=0$.  Its only nonzero
entries are
\begin{equation}\label{eq:cc-matrix}
 \begin{aligned}
 \mathcal B_{ii}&=b_i+Rc_i(w_{i-1}+w_i),
     &b_0=b_m=1,\\[-2pt]
 \mathcal B_{i,i+1}&=Rw_i c_{i+1},
     &b_i=2\quad(0<i<m),\\
 \mathcal B_{i+1,i}&=Rw_i c_i,
     &c_0=c_m=1.
 \end{aligned}
\end{equation}
For $\lambda=50/47$ and $C_{\rm B}=194/3$,
\begin{equation}\label{eq:cc-inverse}
 \|\mathcal B^{-1}\|_{1\to1}\leq1,\qquad
 \sum_{i=0}^m\lambda^{|i-j|}|(\mathcal B^{-1})_{ij}|\leq2,
 \qquad \|\mathcal B^{-1}\|_{\infty\to\infty}\leq C_{\rm B}.
\end{equation}
The constants are independent of $R,m$ and the word.  The algebraic
statement includes simultaneous zero values of interior $c_i$ when
the positive flight lengths are retained.  This extension of the
matrix statement is not a claim that all such arrays are physical.
\end{theorem}
\begin{proof}
For $X_i=O_{i+1}-O_i+R(n_{i+1}-n_i)=\tau_i v_i$,
$d\vartheta_i=v_i^\perp\cdot dX_i/\tau_i$.
At the departure contact $v_i\cdot n_i=c_i$ and at arrival
$v_i\cdot n_{i+1}=-c_{i+1}$.  Differentiating
\eqref{eq:cc-reflection} therefore gives \eqref{eq:cc-matrix}, including
its two boundary rows.  In particular the boundary diagonal is
$1+Rc_iw$, not the interior expression with a spurious second flight.

Conjugate by $S=\operatorname{diag}((-1)^i)$ to obtain
$A=S\mathcal B S=D-C$, where $D$ is the positive diagonal and
$C\geq0$ has zero diagonal.  The excess of the diagonal over the sum
of the off-diagonal absolute values in column $i$ is exactly $b_i$.
Consequently $\|CD^{-1}\|_1<1$ for this finite matrix,
$A^{-1}=D^{-1}\sum_{r\geq0}(CD^{-1})^r\geq0$, and
$\mathbf1^TA\geq\mathbf1^T$ implies
$\mathbf1^TA^{-1}\leq\mathbf1^T$.  This proves the first bound.
No row-dominance assertion is being used.

Fix $j$ and conjugate further by
$W=\operatorname{diag}(\lambda^{|i-j|})$.  Across each edge the
weight ratio is at most $\lambda$.  Since
$a_*=R_+/\tau_-=47/6$ and $a_*(\lambda-1)=1/2$, the column
excess of $WAW^{-1}$ is at least $1-1/2=1/2$ at either endpoint
and at least $2-1=1$ at an interior index.  The same positive-inverse
argument gives its inverse $\ell^1$ norm at most $2$.  As $W_{jj}=1$,
this is the middle assertion in \eqref{eq:cc-inverse}.  In particular
$|(\mathcal B^{-1})_{ij}|\leq2\lambda^{-|i-j|}$.  Sum this bound
over a row, retaining at most two entries at each positive distance:
\[
 \sum_j|(\mathcal B^{-1})_{ij}|
 \leq2\left(1+2\sum_{d\geq1}\lambda^{-d}\right)
 =2\left(1+\frac2{\lambda-1}\right)=\frac{194}{3}.
\]
All steps remain valid when some interior $c_i$ are zero.  The endpoint
case $m=1$ is the same two-by-two calculation, without an interior row.
\end{proof}

\begin{theorem}[A common formal radius domain and physical radius response]
\label{thm:cc-radius}
Every actual regular critical record has a holomorphic continuation
of \eqref{eq:cc-reflection} on a radius disk of size independent of
its collision count and its smallest interior incidence.  More
precisely, put
\begin{equation}\label{eq:cc-analytic-constants}
 \begin{aligned}
 L_R&=4/\tau_-,& K_R&=48/\tau_-+512/\tau_-^2,\\
 s_R&=\min\{1/8,\tau_-/24\},&
 r_R&=\min\{s_R,(8C_{\rm B}K_R)^{-1}\},\\
 \eta_R&=\frac{r_R}{8(1+C_{\rm B}L_R)}.
 \end{aligned}
\end{equation}
For $|\rho|<\eta_R$ the fixed local equations have a unique solution
$\alpha(R+\rho)$ in
$\|\alpha-\alpha(R)\|_\infty\leq r_R/2$, and
\begin{equation}\label{eq:cc-formal-lift}
 \|\alpha(R+\rho)-\alpha(R)\|_\infty
 \leq 2(1+C_{\rm B}L_R)|\rho|.
\end{equation}
This is a formal normal-reflection domain.  On every connected
interval on which its record is actually physical, regular, and
$R\in[R_-,R_+]$, one has the sharper derivative bound
\begin{equation}\label{eq:cc-radius-derivative}
 \max_i|\alpha_i'(R)|\leq D_R,
 \qquad D_R=C_{\rm B}L_R=\frac{38800}{9}.
\end{equation}
Physical continuation through an unvisited tangency or a grazing
collision is not asserted.  A section crossing is not excluded and
its actual selector is reevaluated at the new radius.
\end{theorem}
\begin{proof}
First hold the angles fixed at the reference record.  The formula
$\partial_R\vartheta_i=v_i^\perp\cdot(n_{i+1}-n_i)/\tau_i$
gives $|\partial_R\vartheta_i|\leq2/\tau_-$.  Hence
$\|E_R\|_\infty\leq L_R$.  At every actual regular critical point,
implicit differentiation and Theorem~\ref{thm:cc-full-resolvent} give
$\alpha'=-\mathcal B^{-1}E_R$ and
\eqref{eq:cc-radius-derivative}.

Here are explicit bounds for a domain, rather than just pointwise
implicit differentiation.  Let $z_i$ denote the complex lattice center
$O_i$.  For complex independent increments $(\rho,u)$ use the two
holomorphic flight coordinates
\[
 X_i^\pm=(z_{i+1}^\pm-z_i^\pm)
  +(R+\rho)\bigl(e^{\pm i(\alpha_{i+1}+u_{i+1})}
                         -e^{\pm i(\alpha_i+u_i)}\bigr),
 \qquad z_i^+=z_i,\quad z_i^-=\overline z_i.
\]
The minus coordinate is defined by this expression, not by conjugating
complex variables.  At the real reference point both moduli are
$\tau_i\geq\tau_-$.  In the norm
$\| (\rho,u)\|=|\rho|+\|u\|_\infty$, for increments of norm at
most $s_R$, each flight changes by at most $6s_R\leq\tau_-/4$.
Indeed $|e^{iu}-1|\leq e^{s_R}|u|$, $e^{s_R}<2$, and $R<1/2$.
Thus logarithms continuous from the reference flights are defined
on convex nonzero disks.  On this domain $|R+\rho|<1$, and direct
differentiation gives the conservative bounds
\[
 \|DX_i^\pm\|\leq8,\qquad \|D^2X_i^\pm\|\leq12,
 \qquad |X_i^\pm|\geq\tau_-/2.
\]
The first two inequalities follow by distributing the radius derivative
and at most two adjacent angular derivatives; their operator norms use
the displayed joint norm and do not sum over collision indices.
For $\vartheta_i=(\log X_i^+-\log X_i^-)/(2i)$ up to its fixed
reference constant, the second derivative of the logarithm gives
\[
 \|D^2\vartheta_i\|
 \leq24/\tau_-+256/\tau_-^2,
 \qquad \|D^2E\|\leq K_R.
\]
These are estimates for each component in the supremum norm, uniformly
in the dimension.

Write
$E(R+\rho,\alpha+u)-E(R,\alpha)=\mathcal B u+E_R\rho+N(\rho,u)$.
Taylor's integral formula gives
$\|N\|_\infty\leq(K_R/2)(|\rho|+\|u\|_\infty)^2$ and
$\|D_uN\|\leq K_R(|\rho|+\|u\|_\infty)$.
For $|\rho|\leq\eta_R$ the map
\[
 u\longmapsto-\mathcal B^{-1}(E_R\rho+N(\rho,u))
\]
preserves $\|u\|_\infty\leq r_R/2$: the linear forcing is at most
$r_R/8$ and the nonlinear term at most $r_R/16$.  Its Lipschitz
constant is at most $1/8$.  Iteration from zero gives the unique fixed
point there, holomorphic by locally uniform convergence of the
holomorphic iterates.  Comparing the fixed point with the first
iterate gives
\[
 \|u\|_\infty
 \leq \frac{C_{\rm B}L_R|\rho|+(C_{\rm B}K_R/2)|\rho|^2}{1-1/8}
 \leq 2(1+C_{\rm B}L_R)|\rho|.
\]
This proves the asserted formal domain, also for the zero-incidence
matrix limits with nonzero reference flights and the same reflected
angle identities.  For a regular physical reference, continuity and
its finitely many positive flight/incidence/third-disk margins give
some open physical subinterval; its size need not be uniform.  The
argument has not promoted the formal disk to a common physical radius
disk.  Nor has it frozen a discontinuous exact-record selector.
\end{proof}

\subsection{A collision-uniform derivative of the true critical roof}
Let $J$ be a connected interval of regular physical critical
continuation as in Theorem~\ref{thm:cc-radius}.  Its actual roof is
\[
 a(R)=\sum_{i=0}^{m-1}
  |O_{i+1}-O_i+R(n_{i+1}(R)-n_i(R))|.
\]
Every summand is a real physical flight, not a symbolic roof.

\begin{theorem}[Exact radius variation and uniform critical-clock slope]
\label{thm:cc-clock}
On $J$ one has
\begin{equation}\label{eq:cc-hadamard}
 a'(R)=-c_0-c_m-2\sum_{i=1}^{m-1}c_i
       =-2-2\sum_{i=1}^{m-1}c_i.
\end{equation}
In particular, with $\kappa_0=1/210560$,
\begin{equation}\label{eq:cc-slope}
 2+2\gamma_0\left\lfloor\frac m{94}\right\rfloor
       \leq -a'(R)\leq2m,
 \qquad \kappa_0 m\leq-a'(R)\leq2m.
\end{equation}
For any $R_1<R_2$ in this same physical interval,
\begin{equation}\label{eq:cc-clock-comparison}
 \kappa_0 m(R_2-R_1)
 \leq a(R_1)-a(R_2)\leq2m(R_2-R_1).
\end{equation}
There is also the uniform normalized second-derivative estimate
\begin{equation}\label{eq:cc-second}
 |a''(R)|\leq2(m-1)C_{cR}\leq2mC_{cR},\qquad
 C_{cR}=D_R+\frac{2+2R_+D_R}{\tau_-}.
\end{equation}
Consequently $a/m$ has both a slope separated from zero and a
Lipschitz derivative with constants independent of the collision word.
The same one-sided first-variation identities and bounds hold at a
fixed finite-word nonpenetrating grazing limit of such records.
They do not identify different words across a singularity.
\end{theorem}
\begin{proof}
The full length action $S(R,\alpha)$ has zero angular gradient at a
critical record: the interior equations are specular stationarity,
and the two endpoint derivatives are $-Rp_0$ and $Rp_m$.
Differentiating $S(R,\alpha(R))$ therefore cancels every angular
chain-rule term and leaves
\[
 a'(R)=\partial_RS(R,\alpha)
   =\sum_{i=0}^{m-1}v_i\cdot(n_{i+1}-n_i).
\]
Each summand is $-c_{i+1}-c_i$.  The interior incidences occur exactly
twice; the two normal endpoints occur once and equal one.  This
proves \eqref{eq:cc-hadamard}, including the signs and both boundary
factors.  Corollary~\ref{cor:cb-collision-bending} supplies the first
lower bound in \eqref{eq:cc-slope}.  Since
$\lfloor m/94\rfloor\geq m/94-1$ and $0<\gamma_0<1$,
\[
 2+2\gamma_0\lfloor m/94\rfloor
 \geq (2\gamma_0/94)m,
 \qquad 2\gamma_0/94=1/210560.
\]
The upper bound follows from $0\leq c_i\leq1$.  Integration over
the actual connected physical interval proves
\eqref{eq:cc-clock-comparison}.  In particular no exchange of a
radius supremum and a time or collision limit is used.

For the second derivative, the full response bound gives
$|n_i'|\leq D_R$ and
\[
 |X_i'|\leq2+2R_+D_R,\qquad
 |v_i'|=|(I-v_iv_i^T)X_i'|/\tau_i
       \leq(2+2R_+D_R)/\tau_-.
\]
Differentiate the incoming formula $c_i=-v_{i-1}\cdot n_i$.
It follows that $|c_i'|\leq C_{cR}$, without dividing by $c_i$.
The endpoints remain normal throughout $J$, so their derivatives are
zero.  Differentiating \eqref{eq:cc-hadamard} yields
\eqref{eq:cc-second}.  At a stated finite-word limit, the positive
flight lengths, the uniform reflection inverse and the formal local
continuation supply the one-sided derivatives.  The finite sums of
incidences and the nonpenetrating polygon bound pass to that limit.
No assertion about an arbitrary branch beyond that endpoint follows.
\end{proof}

\begin{corollary}[Branchwise band width and exact radius--roof coarea]
\label{cor:cc-coarea}
On every such connected physical branch, for a physical roof interval
$I_t$,
\begin{equation}\label{eq:cc-band}
 |\{R\in J:a(R)\in I_t\}|
       \leq |I_t|/(\kappa_0m).
\end{equation}
More generally let $\mathscr K_m$ be the full set of regular physical
critical $m$-flight records with the initial disk represented in a
fixed lattice fundamental domain and $R\in[R_-,R_+]$.  Write $x$ for
a record at a given $R$.  Give it any nonnegative Borel weight $f(R,x)$,
which may include the original exact return/displacement/collision
selector and an original optical weight.  For every nonnegative Borel
roof test $\phi$,
\begin{equation}\label{eq:cc-coarea}
 \int_{R_-}^{R_+}\sum_{x:(R,x)\in\mathscr K_m}
       f(R,x)\phi(a_x(R))\,dR
 =\int_{\R}\phi(t)
       \sum_{\substack{(R,x)\in\mathscr K_m\\a_x(R)=t}}
       \frac{f(R,x)}{|a_x'(R)|}\,dt .
\end{equation}
The equality is an identity of positive measures, possibly infinite;
no global integrability is silently assumed.  In every term
\[
 \frac1{2m}\leq\frac1{|a_x'(R)|}\leq\frac1{\kappa_0m}.
\]
Physical records are counted once; coordinate chart overlaps are not
extra records.  The original probability normalization $c_*$ is
unchanged.  For example, inserting a weight with factor
$(4\pi c_*)^{-1}$ leaves that factor exactly once on both sides.
\end{corollary}
\begin{proof}
The derivative in \eqref{eq:cc-slope} has one sign and is separated
from zero, so each connected branch is a $C^1$ strictly monotone
graph over $R$ and a $C^1$ graph over its physical roof.  The ordinary
one-dimensional substitution formula proves \eqref{eq:cc-band} and
\eqref{eq:cc-coarea} on one chart.  To justify the full physical sum,
there are countably many finite lattice words and each finite-dimensional
angle space is second countable.  The invertible full reflection
matrix and the open physical regularity conditions give a countable
cover by charts on open radius intervals.  Assign a physical point
to its first chart in a fixed enumeration.  The assigned subsets
are Borel, and their indicators can be included in $f$ on the
uncut charts.  Substitution for finitely many nonnegative chart terms,
then monotone convergence, gives the displayed equality.  This
allocation deletes representations, never collisions or physical
source points.  The exact selector is evaluated at $(R,x)$, not
transported from another radius or from a tube center.
\end{proof}

\begin{remark}[Fixed-radius shortage is not a radius average]
\label{rem:cc-target-boundary}
The lower clock slope is a new all-word geometric input on the
original radius family.  Formula~\eqref{eq:cc-coarea} is an auxiliary
exact change of variables, not a replacement of the original
probability problem by a radius-averaged limit theorem.  A bound for
crossings summed at different radii does not follow from the
retained momentum kernel at one radius; conversely, a small radius
integral cannot control $\sup_R$.  Neither \eqref{eq:cc-band} nor the
presence of moderately transverse collisions proves the original
$m^{-2}$ roof nonconcentration, the full nongrazing comparison
capacity, or the shortage limit \eqref{eq:v112-remaining}.
Singular/paired endpoint traces, nonguarded critical centers and the
complete noncritical positive source remain part of that limit.
All original microscopic widths remain fixed before the collision
limsup and $B$ remains outermost.  No theorem or proof body of the
inherited raw-LLT and inverse program has been removed.
\end{remark}
